\documentclass[../main.tex]{subfiles}

\begin{document}

\section{Relazioni su un insieme}

\begin{definition}[Relazione]
    Sia $X$ un insieme. Un sottoinsieme $R \subseteq X \times X$ è detto \textbf{relazione su $X$}.
\end{definition}

\begin{definition}[Relazione di equivalenza]
    Una relazione $R \subseteq X \times X$ è detta \textbf{relazione di equivalenza} se soddisfa le seguenti proprietà:
    \begin{enumerate}[label=(\roman*)]
        \item \textbf{riflessità}: $(x,x) \in R$, $\forall x \in X$
        \item \textbf{simmetria}: $(x,y) \in R \implies (y,x) \in R$, $\forall x,y \in X$
        \item \textbf{transitività}: $(x,y) \in R \text{ e } (y,z) \in R \implies (x,z) \in R$, $\forall x,y,z \in X$
    \end{enumerate}
\end{definition}

Se $R$ è una relazione di equivalenza su $X$ e $(x,y) \in R$, scriviamo $x \sim y$, che si legge "$x$ è equivalente a $y$"

\begin{definition}[Classe di equivalenza]
    Sia $x \in X$ e $R$ una relazione di equivalenza su $X$. L'insieme
    \begin{equation*}
        [x] := \{y \in X : x \sim y\}
    \end{equation*}
    è detto \textbf{classe di equivalenza di $x$}. Ogni elemento di $X$ appartiene a esattamente una classe di equivalenza.
\end{definition}

Per non appesantire la notazione, spesso scriveremo $\overline{x}$ al posto di $[x]$.

\begin{definition}[Insieme quoziente]
    L'insieme
    \[
        \nicefrac{X}{\sim} := \{[x] : x \in X\}
    \]
    è detto \textbf{insieme quoziente}. È l'insieme formato da tutte le classi di equivalenza di $X$. L'insieme quoziente rappresenta quindi una "partizione" di $X$ indotta dalla relazione di equivalenza (vedi \ref{prop:1.35}).
\end{definition}

In sintesi, una classe di equivalenza è un sottoinsieme di $X$ contenente elementi equivalenti tra loro, mentre l'insieme quoziente è l'insieme di tutte queste classi.

\begin{definition}[Proiezione canonica]
    La funzione
    \begin{align*}
        \pi : X & \to \nicefrac{X}{\sim} \\
        x       & \mapsto [x]
    \end{align*}
    è detta \textbf{proiezione canonica}.
\end{definition}

\begin{proposition}
    \label{prop:1.35}
    Siano $x,y \in X$. Allora:
    \begin{itemize}
        \item Se $x \sim y$ abbiamo che $[x] = [y]$.
        \item Se $x \nsim y$ abbiamo che $[x] \cap [y] = \varnothing$.
    \end{itemize}
    Quindi:
    \[
        X = \underset{[x] \in \nicefrac{X}{\sim}}{\biguplus} [x]
    \]
    ossia $\nicefrac{X}{\sim}$ è una partizione di X.
\end{proposition}

\begin{example}
    \
    \begin{itemize}
        \item L'uguaglianza "$=$" è una relazione di equivalenza
              su ogni insieme $X$.
        \item Sia $X = \{1,2,\ldots,n\}$. Definiamo si $\mathcal{P}(X)$ la seguente relazione:
              \[
                  A \sim B \iff |A| = |B| \qquad \forall A,B \subseteq X
              \]
              questa è una relazione di equivalenza e $\nicefrac{\mathcal{P}(X)}{\sim}
                  \simeq \{0,1,\ldots,n\}$.

              Se $A \subseteq X$ è tale che $|A| = k \leq n$ allora $|[A]| = \binom{n}{k} := \frac{n!}{k!(n-k)!}$
    \end{itemize}
\end{example}

\begin{definition}[Coefficiente binomiale]
    Il numero $\binom{n}{k}$ è chiamato \textbf{coefficiente binomiale}, questo perché
    \[
        (x+y)^n = \sum_{k=0}^{n} \binom{n}{k} x^n y^{n-k} \qquad \forall x,y \in \mathbb{C}
    \]
\end{definition}

Sia $G$ un gruppo e $H \subseteq G$ un sottogruppo. La relazione $\sim $ su $G$ definita da
\[
    \boxed{
        g_1 \sim g_2 \iff g_1=g_2 h \qquad \text{per qualche } h \in H
    }
\]
è una relazione di equivalenza:
\begin{enumerate}
    \item (Riflessività) $g \sim g : g \cdot e \qquad \forall g \in G, e \in H$
    \item (Simmetria) $g_1 \sim g_2 \implies g_2 \sim g_1 : g_1 = g_2 h \implies g_1 \underbrace{h^{-1}}_{\in H} = g_2$
    \item (Transitività) $g_1 \sim g_2$, $g_2 \sim g_3 \implies g_1 \sim g_3 : g_1 = g_2 h$, $g_2 = g_3 h' \implies g_1 = g_3 \underbrace{h h'}_{\in H} = g_3 h''$, $\forall g_1,g_2,g_3 \in G$
\end{enumerate}
In questo caso l'insieme quoziente lo indichiamo con $\boxed{\nicefrac{G}{H}}$.

\subsection{Insieme quoziente per gruppi abeliani}

Se $G$ è un gruppo abeliano, possiamo definire la seguente operazione "$+$" (commutativa) su
\[
    \nicefrac{G}{H}: [g_1] + [g_2] := [g_1 + g_2],
\]
vediamo che è ben definita:
se $g_{1}' = g_1 + h_1$ e $g_{2}' = g_2 + h_2$ abbiamo $[g_{1}'] = [g_1]$, $[g_{2}'] = [g_2]$ e
\[
    g_{1}' + g_{2}' = g_1 + h_1 + g_2 + h_2 = g_1 + g_2 + h \qquad \text{dove } h = h_1 + h_2 \in H
\]
Quindi
\[
    [g_{1}' + g_{2}'] = [g_1 + g_2]
\]

L'operazione è ovviamente associativa e commutativa, perché lo è quella su $G$.

Inoltre
\[
    [g] + [0] = [g] \qquad \forall [g] \in \nicefrac{G}{H}
\]
dove con "$0$" abbiamo indicato l'identità di $G$. Quindi la classe $[0]$ dell'identità di $G$ è l'identità di $(\nicefrac{G}{H}, +)$.

Infine
\[
    [g] + [-g] = [g-g] = [0] \qquad \forall [g] \in \nicefrac{G}{H}
\]
dove con $-g$ abbiamo indicato l'inverso di $g$ in $G$. Quindi
\[
    -[g] = [-g] \qquad \forall [g] \in \nicefrac{G}{H}
\]
ossia $(\nicefrac{G}{H},+)$ è un gruppo abeliano.

\begin{example}
    \
    \begin{itemize}
        \item Se $H = \{0\} \subseteq G$, allora $\nicefrac{G}{H}$ è isomorfo a $G$. ($\{0\}$ gruppo banale e $G$ gruppo abeliano)
        \item Sia $G = (\mathbb{Z} , +)$, $ n \in \mathbb{N} $. Il sottoinsieme $n\mathbb{Z} := \{nk : k \in \mathbb{Z}\}$ è un sottogruppo di $\mathbb{Z}$.
              \begin{itemize}
                  \item $0 \mathbb{Z} = \{0\}$
                  \item $1 \mathbb{Z} = \{\mathbb{Z} \}$
                  \item $2 \mathbb{Z} = \{\ldots,-4,-2,0,2,4,\ldots\}$
                  \item $3 \mathbb{Z} = \{\ldots,-6,-3,0,3,6,\ldots\}$
                  \item $\ldots$
              \end{itemize}
    \end{itemize}
\end{example}
Definiamo il gruppo abeliano
\[
    \boxed{
        \mathbb{Z}_n := \nicefrac{\mathbb{Z}}{n\mathbb{Z}}
    }
\]
\begin{note}
    $\mathbb{Z}_0 = \nicefrac{\mathbb{Z}}{0 \mathbb{Z}} = \nicefrac{\mathbb{Z}}{\{0\}} \simeq \mathbb{Z}$
\end{note}
Sia $n \geq 0$ e siano $x,y \in \mathbb{Z}$. Allora
\begin{align*}
    x \sim y & \iff x = y + h, \quad h \in n\mathbb{Z}                                                               \\
             & \iff x-y = nk \quad \text{per qualche } k \in \mathbb{Z}                                              \\
             & \iff \text{il resto della divisione di $x$ per $n$ è uguale al resto  della divisione di $y$ per $n$}
\end{align*}

I possibili resti della divisione per $n$ sono $\{0,1,\ldots,n-1\}$.

Quindi
\[
    \mathbb{Z}_n = \{[0],[1],\ldots,[n-1]\} = \{\overline{0}, \overline{1},\ldots,\overline{n - 1}\}
\]
\begin{example}\
    \begin{center}
        \begin{multicols}{2}
            $\mathbb{Z}_2 = \{\overline{0},\overline{1}\} \quad$
            \renewcommand{\arraystretch}{1.25}
            \begin{tabular}{@{\hspace{8pt}}c@{\hspace{8pt}}|@{\hspace{8pt}}c@{\hspace{8pt}}|@{\hspace{8pt}}c@{\hspace{8pt}}}
                +              & $\overline{0}$ & $\overline{1}$ \\ \hline
                $\overline{0}$ & $\overline{0}$ & $\overline{1}$ \\ \hline
                $\overline{1}$ & $\overline{1}$ & $\overline{0}$ \\
            \end{tabular}

            $\overline{1} + \overline{1} = [1 + 1] = [2] = [0]$

            \columnbreak

            $\mathbb{Z}_3 = \{\overline{0},\overline{1},\overline{2}\} \quad$
            \renewcommand{\arraystretch}{1.25}
            \begin{tabular}{@{\hspace{8pt}}c@{\hspace{8pt}}|@{\hspace{8pt}}c@{\hspace{8pt}}|@{\hspace{8pt}}c@{\hspace{8pt}}|@{\hspace{8pt}}c@{\hspace{8pt}}}
                +              & $\overline{0}$ & $\overline{1}$ & $\overline{2}$ \\ \hline
                $\overline{0}$ & $\overline{0}$ & $\overline{1}$ & $\overline{2}$ \\ \hline
                $\overline{1}$ & $\overline{1}$ & $\overline{2}$ & $\overline{0}$ \\ \hline
                $\overline{2}$ & $\overline{2}$ & $\overline{0}$ & $\overline{1}$ \\
            \end{tabular}
        \end{multicols}
    \end{center}
\end{example}

\begin{proposition}
    Sia $G$ un gruppo abeliano e $H \subseteq G$ un sottogruppo. La proiezione canonica
    \[
        \pi : G \rightarrow \nicefrac{G}{H}
    \]
    è un \textbf{morfismo suriettivo di gruppi}
\end{proposition}

Se $G$ è un gruppo finito e $H \subseteq G$ è un sottogruppo, allora
\[
    [g] \in \nicefrac{G}{H} \implies |[g]| = |H|
\]

Infatti $[g] = \{gh : h \in  H\}$ e $gh_1 = gh_2 \implies h_1 = h_2$.

Poiché le classi di equivalenza sono una partizione di $G$, abbiamo
\[
    \abs{G} = \abs{\nicefrac{G}{H}} \cdot \abs{H} \iff \nicefrac{\abs{G}}{\abs{H}} = \abs{\nicefrac{G}{H}}
\]

In particolare la cardinalità o (\textbf{ordine}) di un sottogruppo di un gruppo finito divide la cardinalità del gruppo.
\end{document}